% Generated draft: requires mathematical and source review.
\documentclass[11pt]{article}
\usepackage[margin=25mm]{geometry}
\usepackage{fontspec}
\setmainfont{Latin Modern Roman}
\usepackage{amsmath,amssymb,mathtools,unicode-math}
\setmathfont{Latin Modern Math}
\usepackage{enumitem,needspace,xurl,hyperref}
\setlength{\parindent}{0pt}
\setlength{\parskip}{5pt}
\setlength{\emergencystretch}{3em}
\newcommand{\N}{\mathbb{N}}
\newcommand{\Z}{\mathbb{Z}}
\newcommand{\Q}{\mathbb{Q}}
\newcommand{\R}{\mathbb{R}}
\newcommand{\C}{\mathbb{C}}
\newcommand{\F}{\mathbb{F}}
\newcommand{\A}{\mathbb{A}}
\newcommand{\PP}{\mathbb{P}}
\newcommand{\E}{\mathbb{E}}
\newcommand{\eps}{\varepsilon}
\newcommand{\dd}{\,\mathrm{d}}
\begin{document}
\section*{UnsolvedMath ID: 30001094}
\textbf{Source problem code:} catalogue entry
\subsection{Definitions and mathematical statement}
Let \(G\) be a countable group with the usual discrete-group interpretation. For a function \(\phi:G\to\C\), its Herz--Schur norm \(\|\phi\|_{\mathrm{cb}}\) is the completely bounded norm of the Schur multiplier on \(B(\ell^2G)\) given on matrix entries by
\[
(a_{s,t})\longmapsto(\phi(s^{-1}t)a_{s,t}).
\]
The Cowling--Haagerup constant is
\[
\Lambda(G)=\inf\left\{\sup_i\|\phi_i\|_{\mathrm{cb}}:
\phi_i\text{ finitely supported},\quad
\phi_i(g)\to1\text{ for every }g\in G\right\},
\]
with value infinity when no uniformly bounded approximating family exists. In particular, \(\Lambda(G)=1\) asks for bounds arbitrarily close to one; it does not silently require that the infimum be attained by an approximating family of norm exactly one.

The Haagerup property, equivalently a-T-menability for these groups, means that \(G\) has a proper affine isometric action on a Hilbert space. In displacement terms, for some point \(v\), each set
\[
\{g\in G:\|g\cdot v-v\|\leq R\}
\]
is finite for every finite \(R\geq0\).

The question is the universal implication
\[
\forall G\text{ countable},\qquad
\Lambda(G)=1\quad\Longrightarrow\quad
G\text{ has the Haagerup property}.
\]
The equivalence in the source identifies Haagerup property with a-T-menability; it does not add the reverse implication from that property to \(\Lambda(G)=1\). The assessment separately reports an older Cowling equivalence disproved and leaves this indicated direction as the remaining question.
\subsection{Short English statement}
Must every countable group with Cowling--Haagerup constant one admit a proper affine isometric Hilbert-space action? The source asks this direction, not a new assertion of the historical two-way conjecture.
\subsection{Sources}
\begin{itemize}
\item Local UnsolvedMath registry, permanent ID \(30001094\), source code catalogue entry, ``Cowling--Haagerup Constant and the Haagerup Property.'' No individual proposer is supplied; the assessment refers historically to Cowling's equivalence.
\item \emph{Von Neumann Algebras and Ergodic Theory of Group Actions} (2009), Oberwolfach Reports, DOI 10.4171/owr/2008/49, pp. 2763--2814: \url{https://doi.org/10.4171/owr/2008/49}.
\item Supplied reference, ``Chatterji--Dru\c{t}u--Haglund, counterexample and open converse'': \url{https://www.numdam.org/item/10.1016/j.crma.2007.11.027.pdf}.
\item Supplied literature triage, checked 2026-08-21: status ``open.'' It reports Cowling's historical equivalence disproved but calls the direction \(\Lambda(G)=1\) implies a-T-menability open in its literature assessment. This is supplied dated curation, not independently verified; time-sensitive status and cited claims require independent verification.
\end{itemize}
\end{document}
